\documentclass[11pt,reqno]{amsart}
\usepackage[T1]{fontenc}
\usepackage[utf8]{inputenc}
\usepackage{lmodern}
\usepackage{amsmath,amssymb,amsfonts,amsthm,mathtools,mathrsfs}
\usepackage{enumitem}
\usepackage{booktabs}
\usepackage{microtype}
\usepackage[colorlinks=true,linkcolor=blue,citecolor=blue,urlcolor=blue]{hyperref}
\usepackage[nameinlink,capitalize]{cleveref}
\allowdisplaybreaks
\setlength{\emergencystretch}{2em}

\newtheorem{theorem}{Theorem}[section]
\newtheorem{proposition}[theorem]{Proposition}
\newtheorem{lemma}[theorem]{Lemma}
\newtheorem{corollary}[theorem]{Corollary}
\theoremstyle{definition}
\newtheorem{definition}[theorem]{Definition}
\newtheorem{assumption}[theorem]{Assumption}
\newtheorem{example}[theorem]{Example}
\theoremstyle{remark}
\newtheorem{remark}[theorem]{Remark}

\newcommand{\A}{\mathcal A}
\newcommand{\B}{\mathcal B}
\newcommand{\C}{\mathbb C}
\newcommand{\E}{\mathbb E}
\newcommand{\F}{\mathcal F}
\newcommand{\N}{\mathbb N}
\newcommand{\Pp}{\mathbb P}
\newcommand{\Qp}{\mathbb Q}
\newcommand{\R}{\mathbb R}
\newcommand{\T}{\mathbb T}
\newcommand{\cE}{\mathcal E}
\newcommand{\cP}{\mathcal P}
\newcommand{\Id}{\mathrm{Id}}
\newcommand{\tr}{\operatorname{tr}}
\newcommand{\Var}{\operatorname{Var}}
\newcommand{\Cov}{\operatorname{Cov}}
\newcommand{\Lip}{\operatorname{Lip}}
\newcommand{\osc}{\operatorname{osc}}
\newcommand{\supp}{\operatorname{supp}}
\newcommand{\norm}[1]{\left\lVert #1\right\rVert}
\newcommand{\abs}[1]{\left\lvert #1\right\rvert}
\newcommand{\pair}[2]{\left\langle #1,#2\right\rangle}

\title[Physical-time rough homogenization]{Physical-Time Rough Homogenization for Correlated Moving Open Billiards}
\author{Qian Qi}
\address{Peking University, Beijing 100871, China}
\email{qiqian@pku.edu.cn}
\subjclass[2020]{Primary 60F17, 37A50; Secondary 60H10, 37D50}
\keywords{rough paths, open billiards, physical-time suspension, Green--Kubo covariance, area anomaly}
\date{August 29, 2026}

\begin{document}

\begin{abstract}
We prove a physical-time rough invariance principle on the same moving
specular open billiard constructed in the companion response paper.  A
non-reversible finite-state Markov--Gibbs measure on obstacle itineraries is
pushed to the billiard trapped set.  For an explicit centered collision
observable, the Poisson corrector, Green--Kubo covariance, and antisymmetric
area anomaly are computed exactly:
\[
 \Sigma_{\rm coll}=\frac17I_2,\qquad
 \Gamma_{\rm coll}^{\rm as}
 =-\frac{\sqrt3}{28}
 \begin{pmatrix}0&-1\\1&0\end{pmatrix}.
\]
The area anomaly is nonzero, so the process is neither an independent reset
model nor a reversible surrogate.

The actual nonconstant free-flight roof is retained.  A martingale--coboundary
decomposition gives collision-clock rough convergence, while a Hölder
roof decomposition and monotone inverse-clock theorem give the physical-time
limit.  The physical covariance and area are divided by the mean roof.
The strict geometric response from Paper I therefore produces nonzero
derivatives of both coefficients.  We also prove a random-index window
estimate that is uniform enough for the physical-horizon game in Paper III.

Two exact benchmarks are recorded but not used in the same-platform proof:
the four-branch model supplies exact innovations and algebraic covariance,
while the similarity Lorentz model checks rotation and scaling conventions.
\end{abstract}

\maketitle
\tableofcontents

\section{Introduction}

Response and homogenization are logically distinct.  Paper I differentiates
frozen spectral and roof data.  The present paper proves long-time
probabilistic convergence under one invariant law and then converts collision
count to actual physical time.  All load-bearing objects carry the platform
identifier \texttt{OB3-MG-v1}.

The fast process is deterministic once its initial trapped orbit is chosen.
Probability enters through the specified invariant Markov--Gibbs measure.
The finite symbolic state permits exact Poisson calculations, while the
geometric roof remains the genuine free-flight time of the open billiard.
Thus the model retains correlation, a nonzero second-level anomaly, and a
nonconstant physical clock.

We use martingale approximation \cite{Gordin1969,HallHeyde1980},
rough-path convergence \cite{FrizVictoir2010,FrizHairer2014,
KellyMelbourne2016}, and standard suspension time change.  Every coefficient
needed downstream is displayed explicitly.

\section{The Markov--Gibbs collision process}\label{sec:markov}

Let $(I_k)_{k\in\mathbb Z}$ be the stationary Markov chain on
$\{1,2,3\}$ with
\[
 P=\begin{pmatrix}
 0&2/3&1/3\\
 1/3&0&2/3\\
 2/3&1/3&0
 \end{pmatrix},
 \qquad \pi=(1/3,1/3,1/3).
\]
This is the obstacle-code process of Paper I.  Its law $\nu$ is pushed through
the coding map $\chi_a$ to an invariant probability $\mu_a$ on the billiard
trapped set.

\begin{proposition}[Explicit spectral gap]\label{prop:gap}
The chain is irreducible and aperiodic.  Its nontrivial eigenvalues are
\[
 -\frac12\pm i\frac{\sqrt3}{6},
\]
both of modulus $1/\sqrt3$.
\end{proposition}

\begin{proof}
The rows and columns sum to one, and cycles of lengths two and three are
allowed.  The characteristic polynomial factors as
\[
 (\lambda-1)\left(\lambda^2+\lambda+\frac13\right).
\]
\end{proof}

Define the collision observable
\[
 g_1=(1,0),\qquad
 g_2=\left(-\frac12,\frac{\sqrt3}{2}\right),\qquad
 g_3=\left(-\frac12,-\frac{\sqrt3}{2}\right).
\]
Then $\sum_i\pi_ig_i=0$ and
\[
 C_0:=\sum_i\pi_i g_i\otimes g_i=\frac12I_2.
\]

Let
\[
 J=\begin{pmatrix}0&-1\\1&0\end{pmatrix},\qquad
 A=-\frac12I_2+\frac{\sqrt3}{6}J.
\]

\begin{lemma}[Two-dimensional representation]\label{lem:representation}
For every state $i$,
\[
 \E[g_{I_{k+1}}\mid I_k=i]=Ag_i.
\]
Consequently
\[
 \E[g_{I_k}\mid I_0=i]=A^kg_i.
\]
\end{lemma}

\begin{proof}
Substitute the three rows of $P$ and the displayed vectors.  The second
identity follows by induction.
\end{proof}

\section{Poisson equation and exact coefficients}\label{sec:poisson}

Since $\rho(A)<1$,
\[
 H=(I-A)^{-1}
 =\frac1{14}
 \begin{pmatrix}
 9&-\sqrt3\\
 \sqrt3&9
 \end{pmatrix}.
\]
Set $h_i=Hg_i$.  Then
\[
 (I-P)h=g.
\]
Define
\[
 d_{k+1}=h_{I_{k+1}}-(Ph)_{I_k}.
\]

\begin{proposition}[Martingale--coboundary decomposition]\label{prop:gordin}
The sequence $(d_{k+1})$ is a bounded martingale difference and
\[
 \sum_{k=0}^{n-1}g_{I_k}
 =
 \sum_{k=0}^{n-1}d_{k+1}+h_{I_0}-h_{I_n}.
\]
\end{proposition}

\begin{proof}
Conditional expectation gives
$\E[d_{k+1}\mid I_0,\ldots,I_k]=0$.  The identity is the telescoping form of
$g=h-Ph$.
\end{proof}

Put
\[
 S=A(I-A)^{-1}
 =-\frac5{14}I_2+\frac{\sqrt3}{14}J.
\]

\begin{theorem}[Green--Kubo covariance and area]\label{thm:coefficients}
The collision-clock Green--Kubo covariance is
\[
 \Sigma_{\rm coll}
 =
 C_0+\sum_{k\ge1}
 \left(
 \E[g_{I_0}\otimes g_{I_k}]
 +\E[g_{I_k}\otimes g_{I_0}]
 \right)
 =\frac17I_2.
\]
The canonical piecewise-linear lift has antisymmetric deterministic rate
\[
 \Gamma_{\rm coll}^{\rm as}
 =
 \frac12\sum_{k\ge1}
 \left(
 \E[g_{I_0}\otimes g_{I_k}]
 -\E[g_{I_k}\otimes g_{I_0}]
 \right)
 =-\frac{\sqrt3}{28}J.
\]
\end{theorem}

\begin{proof}
By \cref{lem:representation},
\[
 \E[g_{I_0}\otimes g_{I_k}]
 =C_0(A^k)^T.
\]
Summing the geometric series gives
\[
 \Sigma_{\rm coll}
 =C_0+C_0S^T+SC_0
 =\frac17I_2.
\]
The antisymmetric half-difference is
\[
 \frac12(C_0S^T-SC_0)
 =-\frac{\sqrt3}{28}J.
\]
\end{proof}

\begin{remark}
The covariance is symmetric and positive.  The area tensor is kept separate;
in an RDE it appears in the enhanced driver and may produce a Lie-bracket
drift.  It is never inserted into the diffusion matrix.
\end{remark}

\section{Collision-clock rough convergence}\label{sec:collision-rough}

For $n\ge1$ define
\[
 W_n(t)=n^{-1/2}\sum_{k<nt}g_{I_k}
\]
with polygonal interpolation, and let $\mathbb W_n$ be its canonical
piecewise-linear second level:
\[
 \mathbb W_n(t)
 =
 \frac1n\sum_{0\le j<k<nt}g_{I_j}\otimes g_{I_k}
 +\frac1{2n}\sum_{k<nt}g_{I_k}^{\otimes2}.
\]

\begin{theorem}[Enhanced invariance principle]\label{thm:rough}
For every $p>2$,
\[
 (W_n,\mathbb W_n)
 \Longrightarrow
 \left(B,\mathbb B^{\rm Str}
 +t\Gamma_{\rm coll}^{\rm as}\right)
\]
in the $p$-variation rough-path topology, where $B$ is Brownian motion with
covariance $\Sigma_{\rm coll}$.
\end{theorem}

\begin{proof}
Apply \cref{prop:gordin}.  The bounded martingale differences satisfy
conditional Lindeberg and have ergodic predictable bracket.  The martingale
functional central limit theorem gives the first level.  Discrete
Burkholder inequalities yield, for every $q\ge2$,
\[
 \E\abs{\sum_{k=m}^{n-1}g_{I_k}}^q
 \le C_q\abs{n-m}^{q/2},
\]
and corresponding second-level bounds.  These imply rough tightness.
The symmetric second level is fixed by
\[
 \operatorname{Sym}\mathbb W_{n;s,t}
 =\frac12W_{n;s,t}^{\otimes2}.
\]
The bounded corrector contributes only the stationary antisymmetric average
computed in \cref{thm:coefficients}.  Standard martingale rough-path
identification completes the proof.
\end{proof}

\begin{corollary}[RDE homogenization]\label{cor:rde}
Let $V_0,V_1,V_2$ be $C_b^3$ vector fields.  A consistent Euler recursion
driven by the collision increments converges to the RDE driven by
\[
 \left(B,\mathbb B^{\rm Str}
 +t\Gamma_{\rm coll}^{\rm as}\right).
\]
\end{corollary}

\section{The actual physical roof}\label{sec:roof}

Let $\tau_a:\Sigma\to(0,\infty)$ be the actual free-flight roof from Paper I
and
\[
 \bar\tau_a=\int\tau_a\,d\nu.
\]
It is Hölder and uniformly bounded away from zero.  Define
\[
 S_n^\tau=\sum_{k=0}^{n-1}\tau_a\circ\sigma^k,\qquad
 N_a(t)=\max\{n:S_n^\tau\le t\}.
\]

\begin{lemma}[Clock law and maximal moments]\label{lem:clock}
For every $q\ge2$,
\[
 \E\max_{m\le n}
 \abs{S_m^\tau-m\bar\tau_a}^q
 \le C_qn^{q/2}
\]
uniformly in $a\in I$.  Consequently,
\[
 \sup_{0\le t\le T}
 \abs{\frac1nN_a(nt)-\frac{t}{\bar\tau_a}}
 \longrightarrow0
\]
in probability.
\end{lemma}

\begin{proof}
A two-sided Hölder observable on a mixing finite-type shift is cohomologous to
a one-sided Hölder observable, with bounded transfer coboundary.  The
one-sided transfer operator has a uniform spectral gap.  Solving the Poisson
equation gives a bounded martingale--coboundary decomposition for the centered
roof.  Burkholder and Doob inequalities give the maximal moments.  Since
$\tau_a\ge\tau_*>0$, monotone inversion gives the uniform clock law.
\end{proof}

\begin{theorem}[Physical-time rough limit]\label{thm:physical-rough}
Let
\[
 \widetilde W_n(t)
 =n^{-1/2}\sum_{k<N_a(nt)}g_{I_k},
\]
with its canonical second level.  Then
\[
 (\widetilde W_n,\widetilde{\mathbb W}_n)
 \Longrightarrow
 \left(
 \bar\tau_a^{-1/2}B,\,
 \bar\tau_a^{-1}\mathbb B^{\rm Str}
 +t\bar\tau_a^{-1}\Gamma_{\rm coll}^{\rm as}
 \right).
\]
Equivalently,
\[
 \Sigma_{\rm phys}(a)=\frac1{7\bar\tau_a}I_2,\qquad
 \Gamma_{\rm phys}^{\rm as}(a)
 =-\frac{\sqrt3}{28\bar\tau_a}J.
\]
\end{theorem}

\begin{proof}
Combine \cref{thm:rough,lem:clock} with continuity of monotone time change in
the rough-path topology.  A centered impulse has no additional deterministic
clock correction: the random index differs from its deterministic
approximation by $O_{\Pp}(\sqrt n)$, and a centered partial sum over such a
window is $O_{\Pp}(n^{1/4})$, negligible after $n^{1/2}$ normalization.
\end{proof}

\section{Response of physical coefficients}\label{sec:response}

Paper I proves
\[
 \bar\tau_a'
 =
 -2\int
 \mathbf1_{\{\omega_0=1\}}\cos\varphi_a\,d\nu<0.
\]

\begin{corollary}[Strict coefficient response]\label{cor:response}
The physical covariance and area are differentiable and
\[
 \Sigma_{\rm phys}'(a)
 =-\frac{\bar\tau_a'}{7\bar\tau_a^2}I_2\ne0,
\]
\[
 (\Gamma_{\rm phys}^{\rm as})'(a)
 =\frac{\sqrt3\,\bar\tau_a'}{28\bar\tau_a^2}J\ne0.
\]
\end{corollary}

\begin{proof}
The collision coefficients are independent of $a$ for the chosen fixed
Markov observable.  Differentiate the factors $\bar\tau_a^{-1}$.
\end{proof}

The same formulas arise from the leading eigenvalue of the joint twist
\[
 \mathcal R_{a,q,s}F
 =
 \mathcal R_a(e^{q\cdot g-s\widehat\tau_a}F).
\]
If $\Lambda_a(q)$ solves the zero-pressure equation, then
\[
 D_q^2\Lambda_a(0)=\Sigma_{\rm phys}(a).
\]

\section{A random-index window theorem}\label{sec:random-window}

The physical-horizon game in Paper III requires a quantitative comparison
between a random collision count and a deterministic one.  Put
\[
 n(t)=\left\lfloor\frac{t}{\bar\tau_a}\right\rfloor,\qquad
 R_t=\abs{N_a(t)-n(t)}.
\]

\begin{theorem}[Random-index window]\label{thm:window}
For every $T<\infty$ there is $C_T$ such that, for $t\le T\varepsilon^{-2}$,
\[
 \E R_t\le C_T\sqrt t
\]
and
\[
 \E\abs{
 \varepsilon\sum_{k=n(t)}^{N_a(t)-1}g_{I_k}}
 \le C_T\sqrt\varepsilon,
\]
with the natural reversed sum when $N_a(t)<n(t)$.
\end{theorem}

\begin{proof}
The first estimate follows from \cref{lem:clock} and monotone inversion.
For the second, no stopping-time property is used.  For $q>4$ the forward and
reverse maximal inequalities give
\[
 \E\max_{0\le r\le R}
 \abs{\sum_{k=n}^{n+r-1}g_{I_k}}^q
 \le C_qR^{q/2}.
\]
Split according to $R_t\le C\sqrt t$ and the dyadic events
\[
 2^m\sqrt t<R_t\le2^{m+1}\sqrt t.
\]
The $q$th moment bound for $R_t$ makes the dyadic series summable.  The
unscaled random-window sum is $O(t^{1/4})$ in $L^1$.  With
$t=O(\varepsilon^{-2})$, multiplication by $\varepsilon$ gives
$O(\sqrt\varepsilon)$.
\end{proof}

\section{Benchmark comparison}\label{sec:benchmarks}

The four-branch platform has exact predictable innovations.  Its area anomaly
vanishes unless an additional edge observable is inserted.  It remains the
correct algebraic benchmark for brackets and collision-to-roof normalization.

The similarity Lorentz platform has
\[
 \kappa_a=s_aR_a\kappa_0,\qquad
 \tau_a=s_a\tau_0.
\]
It remains the exact tensorial check for rotations and scales.  Neither
benchmark is imported into \cref{thm:rough,thm:physical-rough}.

\section{Main theorem}

\begin{theorem}[Same-platform physical rough theorem]\label{thm:main}
On \texttt{OB3-MG-v1}, the moving specular billiard, invariant
Markov--Gibbs law, collision covariance, nonzero rough area, actual roof,
physical clock, and geometric coefficient response belong to one dynamical
system.  The data exported to Paper III are
\[
 \Sigma_{\rm coll},\quad
 \Gamma_{\rm coll}^{\rm as},\quad
 \bar\tau_a,\quad
 \Sigma_{\rm phys}(a),\quad
 \Gamma_{\rm phys}^{\rm as}(a),
\]
together with the random-index estimate of \cref{thm:window}.
\end{theorem}

\begin{proof}
Combine \cref{thm:coefficients,thm:rough,thm:physical-rough,cor:response,thm:window}.
\end{proof}

\section*{Scope and review status}

The invariant law is the specified Markov--Gibbs measure on the trapped set,
not an SRB or Liouville measure on a periodic Sinai table.  The manuscript is
prepared for independent external review.

\bibliographystyle{plain}
\bibliography{references}
\end{document}
